\section{问题分析与研究框架}\label{sec:intro}
\subsection{灾害链与研究目标}
吉隆高寒山区地形起伏强烈，冰雪、破碎岩体与沟谷水流共同构成物质和能量传递系统。赛题描述的2026年8月26日灾害具有“冰崩—碎屑流—泥石流—山洪”的链式特征。对于这类过程，气象背景、源区稳定性和下游运动分别对应不同的空间支持与时间尺度：区域观测反映环境演变，结构与水热状态决定源区的抗力和驱动力，释放后的地形、摩阻、孔压及组分控制传播响应。

本文以这一分工组织四问。首先从MET、RADM、SOIL和FLUX提取因素族特征，识别事件前突出的区域环境状态；随后在共同结构几何和参数设计下比较温度、水与损伤机制，建立条件临界关系；再将物理首次过阈、连续变形加速和动态生存三类时间预测模型并列评价；最后通过沿程、二维、孔压和复合相计算刻画释放后的传播场及目标响应。四部分之间交换有物理含义的状态和参数，而不是将某一项统计排序直接当作触发方程。

\subsection{三个时刻与两类空间尺度}
记源区首次机械启动时刻为$t_{\mathrm m}$，进入传播模型的整体释放时刻为$t_{\mathrm r}$，目标$j$达到给定响应阈值的时刻为$t_{\mathrm a,j}$。其关系写为
\begin{equation}
 t_{\mathrm a,j}=t_{\mathrm m}+\Delta_{\mathrm{mr}}+\tau_j,\qquad
 \Delta_{\mathrm{mr}}=t_{\mathrm r}-t_{\mathrm m},\quad
 \tau_j=t_{\mathrm a,j}-t_{\mathrm r}.
 \label{eq:common-clock}
\end{equation}
其中$\Delta_{\mathrm{mr}}$描述机械启动到整体释放的阶段转换，$\tau_j$由给定释放状态和传播模型计算。本文分别保存这三个对象，并以对应的资料与模型评价它们。题面参考时刻用于环境窗口定位，目录时刻及波形资料用于事件阶段分析，目标响应采用释放后的相对时钟。

空间上，气象站作为区域背景支撑，源区和河谷使用米制投影下的几何与地形；对口岸、桥位和设施分别登记代表对象。本文以“区域背景因素—源区条件机制—河谷传播响应”的层级联系这些尺度，使站点异常、物理临界条件和影响足迹具有各自清晰的解释对象。

\subsection{共同实验设计与技术路线}
模型竞争采用共同输入和明确对照。气象因素比较围绕背景年份、天气过程和窗口支持展开；触发机制使用配对参数点，逐项加入温度、水和耦合；时间预测固定发布机会和预测时域；传播模型固定可比地形、物源及边界，并单独检查原生初态、网格与时间步。

\begin{figure}[htbp]
\centering
\begin{tikzpicture}[>=Latex,node distance=0.55cm,box/.style={draw,rounded corners=2pt,align=center,text width=3.45cm,minimum height=1.0cm,font=\small,fill=blue!4},note/.style={font=\footnotesize,align=center}]
\node[box] (d) {多源资料与统一支持\\观测、预报、遥感和地形};
\node[box,right=0.8cm of d] (q1) {因素族识别\\边际、条件贡献、相似过程};
\node[box,below=0.55cm of d] (q2) {结构水热竞争\\稳定裕度、临界关系};
\node[box,right=0.8cm of q2] (q3) {时间预测\\首次过阈、变形、生存};
\node[box,below=0.55cm of q2] (q4) {给定释放传播\\二维、孔压、复合相};
\node[box,right=0.8cm of q4] (out) {联合结果与防灾启示\\因素、时标、过程和目标};
\draw[->] (d)--(q1);
\draw[->] (d)--(q2);
\draw[->] (q1)--(q3);
\draw[->] (q2)--(q3);
\draw[->] (q2)--(q4);
\draw[->] (q3)--(out);
\draw[->] (q4)--(out);
\end{tikzpicture}
\caption{研究技术路线。跨问交换采用同一参数或情景标识；环境因素、机械启动和给定释放后的运动分别建模。}
\label{fig:common-framework}
\end{figure}

本文的主要贡献体现在三个方面。第一，以因素族聚合和多路线竞争减少同一物理因素的重复入选，并同时报告方向与排序稳定性。第二，将机制方程、响应时标、参数补偿和可观测预测放在同一分析中，形成结构—温度—水的条件临界图景。第三，将时间预测和传播实验按信息类型、物理阶段与评价尺度关联，形成可追溯的模型比较框架。由此，复杂模型的价值通过新增过程量和可检验响应呈现。
\AIcredit
